\section{The finite range of the progression theorem}
\label{app:lemmaA-finite}

This appendix completes the proof of \autoref{thm:lemmaA}. It gives the
certificate for the $127$ orders $92\le k\le2730$ with at least three odd
prime factors, the only ones not covered by \autoref{prop:lemmaA-twoprime}
or \autoref{prop:lemmaA-tail}, and proves the two bounds on
Jacobsthal's function used in \autoref{sec:lemmaA-tail} for the orders above
$2730$. The notation is that of \autoref{sec:lemmaA}.

\subsection{The orders up to 2730 with three or more odd primes}

By \autoref{lem:lemmaA-frame}(b), (c) and \autoref{lem:lemmaA-zero}, the
normalised shortest vector of $\Lambda$ at an order $k\le2730$ gives a triple
in $\mathcal V$ with $g\mid2k$ and $3l^4\le16k^2$: indeed
$l^4\le16\cdot2730^2/3$ gives $l^2\le6304<6400$, so $l\le80$. Besides the chord table (C1) of
\autoref{sec:lemmaA-tail}, the finite certificate has two more parts.
\begin{enumerate}
\item[(C2)] \emph{Jacobsthal values.} A list of $224$ pairs $(n,m)$ with
$J(n)\le m$. Each is checked over one period: for every residue $y$ modulo
$n$ some $0\le i<m$ has $\gcd(y+i,n)=1$.
\item[(C3)] \emph{The per-order certificate.} For each of the $127$ even
$k$ with $92\le k\le2730$ that have at least three odd prime factors, either
(C2) contains a pair $(n,m)$ with $n$ the product of the primes of $2k$ such
that $k$ satisfies the tail condition \eqref{eq:lemmaA-tailcond} with bound
$m$ (the \emph{tail alternative}), or, for each $(g;p,q)\in\mathcal V$ with $g\mid2k$ and
$3\bigl(g^2(p^2+q^2)\bigr)^2\le16k^2$, one of the following is verified:
\begin{enumerate}
\item[(i)] a pair $(n,m)$ from (C2) with $n\mid2k$, $\gcd(n,g)=1$, every
prime of $2k$ dividing $ng$, and $2k\psi>mg$, where $\psi$ is the length
of the chord witness in (C1); or
\item[(ii)] for each of the $g^2$ pairs $0\le s,t<g$, with
$(A,E)=(p+\tfrac{2k}{g}s,\,q+\tfrac{2k}{g}t)$: a unit $\nu<100$
satisfying \eqref{eq:lemmaA-good}, or $\gcd(A,E,2k)>1$, or an
$m$ with $0<m<g$ and $2k\mid m(pE-qA)$.
\end{enumerate}
\end{enumerate}
The tail alternative of (C3) holds at $65$ of these $127$ orders, and
alternative (ii) is needed only at $18$ orders, all at most $924$.

\begin{proposition}
\label{prop:lemmaA-finite}
\autoref{thm:lemmaA} holds at the $127$ even orders $92\le k\le2730$ with at
least three odd prime factors.
\end{proposition}

\begin{proof}
If the tail alternative of (C3) holds, then $J(2k)=J(n)\le m$ and $k$
satisfies the tail condition with bound $m$, hence with bound $J(2k)$, and
\autoref{prop:lemmaA-tail} applies. Otherwise, let $v=gu$ be a normalised shortest vector of
$\Lambda$. As noted above, $(g;p,q)\in\mathcal V$, $g\mid2k$ and
$3l^4\le16k^2$, so (C3) applies to it. In case (i), $J(n)\le m$ and $2k\psi>J(n)g$, and
\autoref{prop:lemmaA-table} shows that $(a,e)$ has a witness; the conditions on $n$
are those of \autoref{lem:lemmaA-line}. In case (ii), let $(s,t)$ be the
class given by \autoref{lem:lemmaA-class}. Its two exclusions are ruled out
by that lemma, so the certificate gives a unit $\nu$ showing that $(A,E)$ has a
witness, and then so does $(a,e)$, by \autoref{lem:lemmaA-class}.
\end{proof}


\subsection{Elementary bounds on Jacobsthal's function}
\label{app:lemmaA-jacbounds}

This subsection proves \autoref{cor:lemmaA-jac}, the two inequalities for
Jacobsthal's function used in \autoref{prop:lemmaA-tail}, by a union bound and
a progression trick.

\begin{lemma}
\label{lem:lemmaA-jactools}
\begin{enumerate}
\item[(a)] Suppose every $m$ consecutive integers include at least $c$ that
are prime to $N$, and $\mathcal T$ is a finite set of primes with
$\sum_{p\in\mathcal T}\lceil m/p\rceil<c$. Then $J(n)\le m$ for every $n\ge1$
whose prime factors divide $N$ or lie in $\mathcal T$.
\item[(b)] If $\gcd(f,N)=1$ then $J(fN)\le f\,J(N)$.
\end{enumerate}
\end{lemma}

\begin{proof}
(a) Among $m$ consecutive integers at most $\lceil m/p\rceil$ are multiples
of $p$. So fewer than $c$ of them are divisible by a prime of $\mathcal T$,
and one of the $c$ integers prime to $N$ escapes; it is prime to $n$.

(b) Given a run $y,y+1,\dots,y+fJ(N)-1$, choose $0\le r<f$ with
$y+r\equiv1\pmod f$, and consider $y_i=y+r+fi$ for $0\le i<J(N)$, all in the
run. Each $y_i$ is prime to $f$. With $f'$ an inverse of $f$ modulo $N$,
$f'y_i\equiv f'(y+r)+i\pmod N$ runs through $J(N)$ consecutive integers, so
some $y_i$ is prime to $N$, hence to $fN$.
\end{proof}

Let $\varpi_1<\varpi_2<\dots<\varpi_8$ be the primes $17,19,23,29,31,37,41,43$ and
$\Pi_r=\varpi_1\cdots \varpi_r$ (with $\Pi_0=1$), and recall that $\omega(n)$ is the
number of distinct prime factors of $n$.

\begin{lemma}[at most eight prime factors]
\label{lem:lemmaA-pattern}
Let $n\ge1$ with $\omega(n)\le8$. Let $B$ be the product of the primes
$p\le13$ dividing $n$, and $r$ the number of primes $p>13$ dividing $n$.
Then $B\Pi_r\le n$, and there is an $m$ with $J(n)\le m$ such that
\begin{enumerate}
\item[(i)] if $B$ is even, $178\,m^2\le10\max\bigl(B\Pi_r,\
B(\lfloor2730/B\rfloor+1)\bigr)$;
\item[(ii)] $m\le8$, or $1000\,m<102\,B\Pi_r$.
\end{enumerate}
\end{lemma}

The maximum in (i) is the better of two lower bounds for an order $k$ to
which the lemma will be applied: $k\ge B\Pi_r$, and, as $k$ is a multiple
of $B$ above $2730$, $k\ge B(\lfloor2730/B\rfloor+1)$. So (i) gives the first
inequality of the tail condition, and (ii) gives $J(g)\le8$ or
$J(g)<0.102\,g$, as needed for the band.

\begin{proof}
The $i$-th smallest prime above $13$ is $\varpi_i$, so the primes
$p_1<\dots<p_r$ of $n$ above $13$ satisfy $p_i\ge \varpi_i$. Hence
$n\ge B\,p_1\cdots p_r\ge B\Pi_r$, and
$\sum_i\lceil m/p_i\rceil\le\sum_i\lceil m/\varpi_i\rceil$ for every $m$.

The rest is a finite table. For each of the $64$ sets $\mathcal P$ of primes
$\le13$, with product $B$, and each $r$ with $0\le r\le8-|\mathcal P|$, there is an
entry; this gives $384$ pairs $(B,r)$. For each the table supplies integers
$m$ and $c$ such that
\begin{enumerate}
\item[(P1)] every $m$ consecutive integers include at least $c$ that are
prime to $B$;
\item[(P2)] $\sum_{i=1}^r\lceil m/\varpi_i\rceil<c$;
\item[(P3)] (i) and (ii) hold for this $m$.
\end{enumerate}
Condition (P1) is checked by sliding a window of length $m$ over one period
$\{0,\dots,B-1\}$ of the starting point; (P2) and (P3) are arithmetic. The
largest $m$ needed is $84$. Given the table, (P1), (P2)
and the inequality above give exactly the situation of
\autoref{lem:lemmaA-jactools}(a), with $N=B$ and $\mathcal T$ the primes of
$n$ above $13$, so $J(n)\le m$.
\end{proof}

\begin{lemma}[at least nine prime factors]
\label{lem:lemmaA-crude}
Let $n$ have $\omega=\omega(n)\ge9$ distinct prime factors. Let $B$ be the product of the primes $p\le \omega$
dividing $n$, and $R$ the product of the remaining $r$ primes of $n$. Then
$BR\le n$, $J(n)\le B(r+1)$, and
\[
178\,B(r+1)^2\le10R,\qquad1000(r+1)<102R.
\]
\end{lemma}

\begin{proof}
Every prime of $R$ exceeds $\omega\ge r$, so $r+1$ consecutive integers include
at most one multiple of each, and \autoref{lem:lemmaA-jactools}(a) (with
$N=1$, $c=m=r+1$) gives $J(R)\le r+1$. Since $\gcd(B,R)=1$,
\autoref{lem:lemmaA-jactools}(b) gives $J(BR)\le B(r+1)$, and $J(n)=J(BR)$.

For the inequalities, the $r$ primes of $R$ are distinct integers above
$\omega$, so $R\ge(\omega+1)(\omega+2)\cdots(\omega+r)$. The other $\omega-r$ primes are at most
$\omega$, so $B$ divides the primorial $\omega\#=\prod_{p\le \omega}p$ and $r\ge \omega-\pi(\omega)$. If
$\omega\le63$, it suffices to check that
$178\,(\omega\#)(r+1)^2\le10\prod_{i=1}^r(\omega+i)$ and
$1000(r+1)<102\prod_{i=1}^r(\omega+i)$ for every $9\le \omega\le63$ and
$\omega-\pi(\omega)\le r\le \omega$, a finite computation. If $\omega\ge64$, use the classical bound $\omega\#\le4^\omega$ and the
count $\pi(\omega)\le \omega/2+1$ (the odd primes up to $\omega$ are distinct odd numbers
in $[3,\omega]$). The bound on $\omega\#$ holds for $\omega\le2$ by inspection and follows
by induction: for even $\omega>2$, $\omega\#=(\omega-1)\#$; for $\omega=2m+1$, the primes in
$(m+1,2m+1]$ all divide $\binom{2m+1}{m}$, which is at most $4^m$ because it
and $\binom{2m+1}{m+1}$ are equal and sum to at most $2^{2m+1}$, so
$(2m+1)\#\le(m+1)\#\,4^m\le4^{m+1}4^m$. Then $r\ge \omega/2-1\ge31$,
$\omega\le2r+2$, $\prod_{i=1}^r(\omega+i)\ge64^r$ and $B\le4^\omega\le16\cdot16^r$, so
$178B(r+1)^2\le16^r\cdot2848(r+1)^2\le16^r\cdot10\cdot4^r=10\cdot64^r$,
using $2848(r+1)^2\le10\cdot4^r$ for $r\ge8$ (true at $r=8$, and the right
side grows faster). The second inequality is then immediate.
\end{proof}

\begin{corollary}
\label{cor:lemmaA-jac}
\begin{enumerate}
\item[(a)] Every even $k>2730$ satisfies $178\,J(2k)^2\le10k$, that is,
$17.8\,J(2k)^2\le k$.
\item[(b)] Every $g\ge1$ satisfies $J(g)\le8$ or $1000\,J(g)<102\,g$.
\end{enumerate}
\end{corollary}

\begin{proof}
(a) Note $J(2k)=J(k)$. If $\omega(k)\le8$, apply
\autoref{lem:lemmaA-pattern} to $n=k$. Then $B$ is even, and $k\ge B\Pi_r$;
also $k$ is a multiple of $B$ exceeding $2730$, so
$k\ge B(\lfloor2730/B\rfloor+1)$. So $178J(k)^2\le178m^2\le10k$. If
$\omega(k)\ge9$, \autoref{lem:lemmaA-crude} gives
$178J(k)^2\le178B^2(r+1)^2\le10BR\le10k$.
(b) If $\omega(g)\le8$, \autoref{lem:lemmaA-pattern}(ii) gives $J(g)\le m\le8$
or $1000J(g)\le1000m<102B\Pi_r\le102g$. If $\omega(g)\ge9$,
\autoref{lem:lemmaA-crude} gives $1000J(g)\le1000B(r+1)<102BR\le102g$.
\end{proof}

